% =============================================================================
% Kapitel 9: Fazit und Ausblick
% ARBEITSGERUEST / REVIEW CANDIDATE
% Stand: 03.09.2026
%
% Dieses Kapitel bleibt bewusst kompakt.
%
% Kapitel 8:
% - interpretiert Ergebnisse
% - diskutiert Literaturbezug
% - behandelt Limitationen
% - beantwortet HF / TF1 / TF2 / TF3
%
% Kapitel 9:
% - zieht die Schlusslinie
% - benennt den Beitrag der Arbeit
% - beschreibt belastbar abgegrenzte Weiterentwicklungspfade
%
% NICHT:
% - Forschungsfragen erneut ausfuehrlich beantworten
% - neue Ergebnisse einfuehren
% - neue Literaturdiskussion aufmachen
% - offene technische Backlog-Punkte unstrukturiert aufzaehlen
%
% =============================================================================

\section{Fazit und Ausblick}
\label{sec:conclusion}

% =============================================================================
\subsection{Zusammenfassung der Arbeit}
\label{sec:conclusion_summary}
% =============================================================================

% ZIEL:
% Ca. 1 bis 1,5 Seiten.
%
% Den roten Faden der Arbeit in stark verdichteter Form zusammenfassen.
%
% EMPFOHLENE REIHENFOLGE:
%
% 1. Ausgangsproblem:
%    heterogene Legacy-Engineering-Informationen
%    -> fehlende direkte, belastbare Ueberfuehrung in SysML v2
%
% 2. Ziel:
%    Konzeption einer Informationsverarbeitungsarchitektur
%    fuer kontrollierte KI-gestuetzte Bewertung, Synthese und Modellbildung.
%
% 3. Vorgehen:
%    Literatur
%    -> Architekturkonzeption
%    -> modellbasierte Beschreibung
%    -> prototypische Umsetzung
%    -> formative Verifikation
%    -> finale E2E-/Multi-Source-/SYSIDE-/Human-Validation
%
% 4. finale Architektur:
%
% Project / Source Context
% -> source-grounded Processing
% -> Human Engineering Review
% -> Approved Engineering Information
% -> Project Fit / Multi-Source
% -> Model Derivation / Placement
% -> Target-Model Formulation
% -> deterministic Assembly
% -> deterministic SysML-v2 Generation
% -> technical Validation
% -> Final Human Review
% -> Publication
%
% 5. wichtigste Resultate:
%
% - Single-Source E2E erfolgreich
% - genuine Multi-Source fuer Project 308131 erfolgreich
% - externe SYSIDE-Validierung erfolgreich
% - Human Release / immutable Publication erreicht
% - finale Regression 6335 passed / 12 skipped
%
% WICHTIG:
% Diese Punkte nur knapp nennen.
% Interpretation bereits in Kapitel 8.


% =============================================================================
\subsection{Beitrag der Arbeit}
\label{sec:conclusion_contribution}
% =============================================================================

% Ziel:
% klar sagen, was die Arbeit ueber die reine Implementierung hinaus leistet.
%
% Nicht zehn kleine Beitraege nennen.
% Besser 3 bis 4 belastbare Hauptbeitraege.


\subsubsection{Konzeptioneller Beitrag}
\label{sec:conclusion_conceptual_contribution}

% KERNGEDANKE:
%
% Entwicklung einer responsibility-basierten
% Informationsverarbeitungsarchitektur fuer die kontrollierte
% Ueberfuehrung heterogener Engineering-Informationen in SysML v2.
%
% Zentrale Prinzipien:
%
% - Source Grounding und Provenance
% - shared Subject Identity vor Persona Interpretation
% - explizite Human Engineering Authority
% - Trennung von Engineering Meaning und Target Representation
% - explizite Model Placement / Target-Model Formulation
% - deterministic Assembly / Serialization
% - separate technical Validation und Human Approval
%
% NICHT:
% als universelle Referenzarchitektur bezeichnen.
%
% Besser:
% Architekturkonzept fuer den untersuchten Problemraum.


\subsubsection{Methodischer Beitrag}
\label{sec:conclusion_methodological_contribution}

% KERNGEDANKE:
%
% Die Architektur wurde nicht nur entworfen,
% sondern durch einen iterativen formativen Zyklus weiterentwickelt:
%
% Architecture Assumption
% -> Implementation Slice
% -> Focused Verification
% -> Finding
% -> Root Cause
% -> Architecture / Contract Correction
% -> Retest
% -> Regression
% -> Acceptance
%
% Wert:
% Findings wurden als Evidenz erhalten,
% nicht durch Neustart / Glattziehen aus der Historie entfernt.
%
% WICHTIG:
% ADRs = Design Rationale / project evidence.
% Keine Methodik behaupten, die wissenschaftlich nicht sauber belegt wurde.


\subsubsection{Prototypischer und technischer Beitrag}
\label{sec:conclusion_technical_contribution}

% KERNGEDANKE:
%
% Der entwickelte PoC realisiert den thesis-relevanten Informationspfad
% technisch end-to-end.
%
% Belegt durch:
%
% - source-local Processing
% - Human Review
% - AEI
% - Project Fit
% - Model Candidate / Placement
% - IEM
% - SysML-v2 Generation
% - SYSIDE Validation
% - Final Human Review / Publication
% - Genuine Multi-Source Retest
%
% NICHT:
% Production Readiness behaupten.


\subsubsection{Beitrag zur Human-Authority- und Assurance-Frage}
\label{sec:conclusion_human_authority_contribution}

% KERNGEDANKE:
%
% Die Arbeit operationalisiert Human in the Loop nicht als einen
% pauschalen Approve-Schritt, sondern als gestufte Engineering Authority.
%
% Relevante Entscheidungsebenen:
%
% - Engineering Meaning
% - Ambiguity / Variance
% - Model Placement
% - Target-Model Formulation / Quality
% - Final Model Review
% - Publication
%
% Gleichzeitig:
% Interaktion soll auf materielle Entscheidungen konzentriert werden.
%
% Diese Aussage kurz halten.
% Die eigentliche Diskussion steht in Kapitel 8.


% =============================================================================
\subsection{Ausblick}
\label{sec:conclusion_outlook}
% =============================================================================

% AUSDRUECKLICHE TRENNUNG:
%
% - implemented today
% - conceptually enabled
% - future work
%
% Diese Unterscheidung insbesondere beim Adaptive Learning durchhalten.


\subsubsection{Erweiterung auf weitere Engineering-Aufgaben}
\label{sec:outlook_reuse}

% ZENTRALER AUSBLCKSPUNKT:
%
% Die Architektur ist nicht ausschliesslich fuer genau eine einzelne
% Legacy-to-SysML-v2-Aufgabe interessant.
%
% Durch:
% - modulare Processing Boundaries
% - task-spezifische LLM Calls
% - Evidence / Provenance
% - Human Review
% - deterministic downstream responsibilities
%
% koennen zukuenftig weitere klar abgegrenzte Engineering Tasks
% ueber dieselbe Governance- und Authority-Infrastruktur laufen.
%
% WICHTIG:
% Nicht behaupten, dass andere Tasks bereits trainiert / validiert wurden.
%
% Formulierungslogik:
%
% Die Architektur legt eine Wiederverwendung fuer weitere Engineering-Aufgaben nahe.
% Die konkrete Uebertragbarkeit muss fuer jede Aufgabe separat validiert werden.
%
% MOEGLICHE BEISPIELKATEGORIEN, NUR WENN SPATER GEWOLLT:
% - weitere Information-Extraction-Aufgaben
% - Classification / Mapping Tasks
% - Modellqualitaetspruefung
% - Change Impact Assessment
% - Requirements / Architecture transformation support
%
% Keine lange Fantasieliste.
%
% PROJEKTEVIDENZ:
% Thesis_Outlook_Adaptive_Human_Feedback_Learning.md
% -> "reuse/portability applications of the Turing Generator"


\subsubsection{Adaptive Human-Feedback-Lernarchitektur}
\label{sec:outlook_adaptive_learning}

% PROJEKTSTATUS:
% Future concept / thesis outlook
% NICHT Teil der implementierten v0.3.0-Baseline.
%
% PROJEKTDOKUMENT:
% collaboration/checkpoints/Thesis_Outlook_Adaptive_Human_Feedback_Learning.md
%
% KERNIDEE:
%
% Persistierte Human-in-the-Loop-Entscheidungen werden nicht nur
% als Engineering Authority genutzt,
% sondern koennen zukuenftig als strukturiertes supervisiertes Feedback
% fuer weitere LLM-gestuetzte Engineering Tasks dienen.
%
% KEIN:
% uncontrolled online self-modification.
%
% ZIEL:
% controlled continual learning.


\paragraph{Konzeptioneller Lernpfad}
\label{sec:outlook_adaptive_learning_loop}

% ZUKUENFTIGER FLOW:
%
% LLM Proposal
% -> Human Review
% -> Accept as generated / Modify / Reject / Defer
% -> persist exact delta + provenance
% -> detect recurring correction patterns
% -> train / update candidate learned policy
% -> offline / shadow evaluation
% -> governed promotion decision
% -> version learned policy
% -> Adaptive Inference Gateway
% -> next inference cycle
%
% WICHTIG:
% gelernter Zustand ist ein versioniertes / pruefbares Artefakt.
% Keine stille Aenderung produktiver Engineering Authority.


\paragraph{First-Pass Human Acceptance Rate als moegliche Messgroesse}
\label{sec:outlook_acceptance_metric}

% PRIMARY DESCRIPTIVE METRIC:
%
% First-Pass Human Acceptance Rate
%
% accepted without Human content change
% -------------------------------------
% all Human-decided proposals
%
% In LaTeX spaeter ggf. Formel.
%
% Messung bevorzugt ueber:
% exact machine proposal content / fingerprint
% vs.
% Human-reviewed result
%
% NICHT allein ueber UI Outcome Label.
%
% ZUSAETZLICHE HUMAN INTERVENTION METRIKEN:
%
% - accepted as generated
% - accepted with modification
% - rejected
% - deferred / out of scope
% - changed fields
% - type / magnitude of change
% - recurring correction patterns
% - review effort, soweit messbar
%
% WICHTIG:
% Acceptance Rate ist KEIN alleiniger Optimierungs-Score.


\paragraph{Guardrails fuer adaptives Lernen}
\label{sec:outlook_adaptive_guardrails}

% Optimierungslogik:
%
% maximize:
% - first-pass Human acceptance
% - reduction of Human correction effort
%
% subject to:
% - source grounding preserved
% - coverage sufficient
% - traceability complete
% - technical validation passed
% - Human authority rules satisfied
% - uncertainty remains visible
% - no trivial / overconservative proposals only to increase acceptance
%
% WICHTIGER THESIS-PUNKT:
%
% Eine steigende Acceptance Rate darf NICHT automatisch
% Human Approval ersetzen.


\paragraph{Nutzung wiederkehrender Human-Korrekturen}
\label{sec:outlook_tacit_knowledge}

% LANGFRISTIG BESONDERS INTERESSANT:
%
% Nicht nur Accept/Reject lernen,
% sondern wiederkehrende inhaltliche Aenderungsmuster.
%
% PROJEKTBEISPIEL:
%
% Machine repeatedly infers purpose from capability statement
% -> Engineer repeatedly restores source-supported capability semantics
% -> recurring correction detected
% -> candidate learned rule:
%    do not infer purpose when source supports only capability
%
% INTERPRETATION FUER DEN AUSBLICK:
%
% Human Review koennte dadurch neben Quality Gate
% zu einer strukturierten Quelle wiederverwendbaren Expertenwissens werden.
%
% VORSICHT:
% Nicht behaupten, dass tacit knowledge vollstaendig extrahierbar ist.
% Nur "Teile wiederkehrender Expertenkorrekturen" / "Muster".


\paragraph{Moeglicher Reifungspfad}
\label{sec:outlook_learning_maturity}

% REPO-KONZEPT:
%
% Stage 1:
% 100% Human Review
%
% Stage 2:
% learned assistance reduces correction effort
%
% Stage 3:
% risk-/confidence-based Human review focuses on exceptions
%
% Stage 4:
% bounded classes may become candidates for automatic approval
%
% EXTREM WICHTIG:
%
% Stage 4 benoetigt eine neue explizite
% Governance-/Requirements-Entscheidung.
%
% Es ist KEINE automatische Folge hoher Acceptance Rates.


\subsubsection{Erweiterung der semantischen und modellseitigen Abdeckung}
\label{sec:outlook_target_coverage}

% ZIEL:
% die wichtigsten echten technischen Weiterentwicklungen zusammenfassen.
%
% KANDIDATEN:
%
% - breitere SysML-v2-Construct-Coverage
% - generalisierte Target-Model Formulation
% - weitere Relationship families
% - verbesserte predicate-variant consolidation
% - weitere Source Types
%
% RELEVANTE FINDINGS:
% SEM-009
% SEM-011
% SEM-015
%
% NICHT:
% komplette Finding-/Backlog-Liste wiederholen.


\subsubsection{Weiterentwicklung der Multi-Source-Verarbeitung}
\label{sec:outlook_multisource}

% STATUS HEUTE:
% thesis MVP endet bewusst bei:
%
% source admissibility / Project Fit
% + provenance-preserving downstream flow.
%
% ZUKUENFTIG:
%
% project-level semantic reconciliation
% change control
% equivalent / complementary / conflicting / superseding information
% source revisions
% explicit Project Engineering Authority
%
% WICHTIG:
% kein synthetic merged AEI.
%
% source provenance darf semantic grouping nicht automatisch bestimmen.
%
% PROJEKTEVIDENZ:
% Zusatz im Adaptive Human Feedback Learning checkpoint
% nach BLK-002 closeout.


\subsubsection{Betriebs- und Skalierungsfaehigkeit}
\label{sec:outlook_scalability}

% Nur kompakt:
%
% - Background Processing
% - Cancellation / recovery
% - Reprocessing / supersession
% - Runtime / Token / Cost optimization
% - groessere Datensaetze
% - weitere Source Types
%
% Nicht zu viel Raum geben.
% Das sind technische Weiterentwicklungen, nicht Hauptbeitrag.


\subsubsection{Weiterfuehrende wissenschaftliche und industrielle Validierung}
\label{sec:outlook_validation}

% ZUKUENFTIGE UNTERSUCHUNGEN:
%
% - unabhaengige Evaluation
% - groessere / realere Brownfield-Datensaetze
% - weitere Engineering Domains
% - Vergleich unterschiedlicher LLM- / Agentenstrategien
% - quantitative Human-effort measurement
% - quantitative acceptance / intervention measurement
% - langfristige Studie zum Adaptive Human Feedback Learning
%
% Safe Demo:
% je nach Ergebnis als erster externer Praxisimpuls,
% aber NICHT als Ersatz fuer groessere Validierung.
%
% MOEGLICHE LANGFRISTIGE FRAGEN:
%
% - Wie stabil ist First-Pass Acceptance ueber Projekte / Domains?
% - Welche Korrekturmuster sind tatsaechlich wiederverwendbar?
% - Wann reduziert Lernen Review-Aufwand, ohne Coverage / Grounding zu verschlechtern?
% - Wie muss Shadow Evaluation fuer learned policies gestaltet sein?
% - Welche Governance ist fuer risk-based review erforderlich?


% =============================================================================
\subsection{Abschliessende Einordnung}
\label{sec:conclusion_closing}
% =============================================================================

% Optionaler letzter kurzer Abschnitt / letzter Absatz.
%
% Ziel:
% Arbeit mit einer klaren, aber nicht werblichen Schlussaussage beenden.
%
% Moegliche inhaltliche Linie:
%
% Die Arbeit zeigt im abgegrenzten PoC,
% dass KI-gestuetzte Legacy-to-SysML-v2-Verarbeitung nicht allein
% durch leistungsfaehige LLMs bestimmt wird,
% sondern wesentlich durch die Architektur der Informations-, Authority-,
% Provenance- und Validierungsgrenzen.
%
% Der entwickelte Prototyp ist damit weniger als fertiges Produkt zu verstehen,
% sondern als technisch erprobte Instanziierung eines Architekturkonzepts.
%
% Der Ausblick auf weitere Engineering Tasks und
% kontrolliertes Human-Feedback Learning knuepft direkt an
% die heute bereits vorhandenen modularen und versionierten Boundaries an.
%
% WICHTIG:
% Kein neuer Claim.
% Kein Marketing-Satz.
% Keine grossen Zukunftsversprechen.


% =============================================================================
% EMPFOHLENE ABBILDUNGEN / TABELLEN
% =============================================================================
%
% PRIORITAET 1
% Abb. 9.1 Adaptive Human-Feedback Learning Loop
%
% Human Review
% -> intervention delta
% -> learning
% -> shadow validation
% -> governed promotion
% -> Adaptive Inference Gateway
% -> future proposal
%
% Sehr geeignet fuer Thesis + Kolloquium.
%
% PRIORITAET 2
% Tab. 9.1 Implementiert / konzeptionell ermoeglicht / Future Work
%
% Beispiel:
%
% Persisted Human Decisions:
% implemented
%
% First-Pass Acceptance measurement:
% conceptually enabled / future measurement
%
% Learned Policy Training:
% future work
%
% Adaptive Inference Gateway:
% future work
%
% Risk-based Review:
% future work
%
% Automatic Approval bounded classes:
% long-term future, new governance decision required
%
% Diese Tabelle verhindert Overclaiming.


% =============================================================================
% LITERATURBEDARF KAPITEL 9
% =============================================================================
%
% Kapitel 9 sollte keine neue grosse Literaturbasis eroeffnen.
%
% Fuer den Adaptive-Learning-Ausblick kann spaeter,
% wenn passend und bereits in Kapitel 8 diskutiert:
%
% - Human Feedback Learning
% - continual learning
% - active learning
% - human-in-the-loop adaptation
%
% kurz referenziert werden.
%
% Grundsatz:
% Literaturargumentation primaer Kapitel 8.
% Kapitel 9 verweist knapp auf dort etablierte Erkenntnisse.


% =============================================================================
% PROJEKTEVIDENZ FUER KAPITEL 9
% =============================================================================
%
% FINAL STATE:
% collaboration/checkpoints/
% 2026-09-02_catia_modeling_complete_thesis_writing_transition_ssot.md
%
% ADAPTIVE LEARNING OUTLOOK:
% collaboration/checkpoints/
% Thesis_Outlook_Adaptive_Human_Feedback_Learning.md
%
% FINDINGS / OPEN COVERAGE:
% collaboration/audits/wp12_findings.md
%
% MULTI-SOURCE OUTLOOK:
% BLK-002 closeout content in adaptive learning checkpoint
%
% WICHTIG:
% Future concept bleibt klar vom implementierten Baseline-Stand getrennt.


% =============================================================================
% OFFENE ARBEIT VOR AUSFORMULIERUNG
% =============================================================================
%
% [ ] Kapitel 8 finalisieren, bevor Kapitel 9 geschrieben wird.
%
% [ ] Exakte finalen RQ-Antworten aus Kapitel 8 als Basis fuer 9.1 verwenden.
%
% [ ] Safe Demo Ergebnis beruecksichtigen.
%
% [ ] First-Pass Acceptance Rate nur als Future Metric darstellen,
%     sofern fuer aktuelle Runs keine belastbare retrospektive Messung vorliegt.
%
% [ ] Adaptive-Learning-Diagramm erstellen.
%
% [ ] Implementiert / enabled / future Tabelle erstellen.
%
% [ ] Ausblick nicht mit OBS-/Backlog-Liste aufblasen.


% =============================================================================
% CLAIM-CHECK VOR DEM SCHREIBEN
% =============================================================================
%
% [ ] Keine Forschungsfragen erneut ausfuehrlich beantwortet?
% [ ] Keine neue Empirie?
% [ ] Beitrag klar von Resultaten getrennt?
% [ ] Keine Production-Readiness-Behauptung?
% [ ] Reuse fuer weitere Tasks als Potenzial, nicht als validierter Fakt?
% [ ] Adaptive Learning klar Future Work?
% [ ] Kein uncontrolled online learning suggeriert?
% [ ] Acceptance Rate nicht als alleinige Qualitaetsmetrik?
% [ ] Human Authority bleibt erhalten?
% [ ] Automatic Approval nur als langfristige, neue Governance-Entscheidung?
% [ ] Tacit Knowledge nur vorsichtig als wiederkehrende Korrekturmuster?
% [ ] Multi-Source Future Scope korrekt abgegrenzt?
% [ ] Keine lange technische Backlog-Liste?
% [ ] Kapitel insgesamt kompakt?
% [ ] Keine Gedankenstriche als Satzzeichen im finalen Fliesstext?
% [ ] Kein unpersoenliches "man"?
% [ ] Tabellenueberschrift oben, Abbildungsunterschrift unten?
%
% =============================================================================
